% appendix-numerical-combinatorics.tex — \input'd from appendix-numerical.tex
%
% Exact rational polytope data and transition feasibility.
%
% Sources: agent-written from geom/rational.rs code and ablation experiment
%   (lem:transition-feasibility).
%
% Structure:
%   - Intro: rational coordinates, exact computation
%   - Remark: exact quantities from rational halfspaces
%   - Lemma: transition feasibility (restated from ablation, self-contained)
%   - Remark: numerical pruning via three-valued predicates
%   - Definition: simple polytope
%   - Definition: numerically robust simple polytope

% [TODO: JÖRN - Full rewrite of A.2 per discussion.
%  Key changes from previous version:
%  (1) No "combinatorial data" wrapper — just the quantities directly.
%  (2) Transition feasibility lemma restated here (duplicate of ablation.tex,
%      canonical location TBD during final polish).
%  (3) Simple polytope defined later, not assumed upfront.
%  (4) Numerically robust = singular value condition on vertex matrices.
%  Please review: lemma statement, proof, and the two definitions.]

The algorithms in Section~\ref{sec:algorithms} operate on a polytope
$K = K(n, h) = \{x \in \mathbb{R}^4 : \langle n_i, x \rangle \leq h_i,
\; i = 1, \ldots, F\}$.
Our input polytopes have rational halfspace data:
normals $n_i \in \mathbb{Q}^4 \setminus \{0\}$
(not necessarily unit vectors, unlike the unit normals
in Section~\ref{sec:algorithms})
and heights $h_i \in \mathbb{Q}_{>0}$.
This allows several combinatorial quantities to be
computed exactly over~$\mathbb{Q}$, eliminating
numerical predicates from all discrete decisions.

\begin{remark}[Exact quantities from rational coordinates]%
\label{rem:exact-quantities}
Given a polytope $K = K(n, h)$ with rational halfspace data,
the following quantities are computed exactly:
\begin{enumerate}
\item \emph{Vertices.}
  For each 4-element subset $S \subset \{1, \ldots, F\}$
  with $\det(N_S) \neq 0$ (where $N_S$ is the $4 \times 4$
  matrix with rows~$n_i^\top$, $i \in S$),
  the candidate vertex $v_S = N_S^{-1} h_S$ is computed
  via Cramer's rule over~$\mathbb{Q}$.
  The subset~$S$ defines a vertex of~$K$ if
  $\langle n_i, v_S \rangle \leq h_i$ for all $i \notin S$.
\item \emph{Vertex--facet incidence.}
  The incidence matrix $E \in \{0, 1\}^{V \times F}$
  has $E_{j,i} = 1$ if and only if vertex~$v_j$
  lies on facet~$F_i$, i.e.\ $\langle n_i, v_j \rangle = h_i$.
  % Downstream: Rust stores this as vertex descriptors S_j = {i : E_{j,i} = 1}.
\item \emph{Symplectic sign matrix.}
  For each pair of facets $i, k$ with $1 \leq i < k \leq F$,
  \[
    D(i, k) = \operatorname{sign}\bigl(
    \omega_0(n_i, n_k)\bigr) \in \{+, -, 0\},
  \]
  where $\omega_0(u, v) = u_1 v_3 - u_3 v_1
  + u_2 v_4 - u_4 v_2$ is the standard symplectic form
  in coordinates $(q_1, q_2, p_1, p_2)$.
  Since $n_i \in \mathbb{Q}^4$, the value
  $\omega_0(n_i, n_k) \in \mathbb{Q}$ is exact and
  its sign is determined without numerical error.
  % Downstream: Rust computes D(i,k) only for vertex-adjacent pairs
  % as an optimization. Non-adjacent pairs are pruned earlier.
\end{enumerate}
\end{remark}

\begin{lemma}[Transition feasibility]%
\label{lem:transition-feasibility-appendix}
% [TODO: JÖRN - Restated from ablation experiment
%  (Lemma~\ref{lem:transition-feasibility}).
%  Canonical location TBD during final polish.]
Let $K \subset \mathbb{R}^4$ be a convex polytope with
facets $F_1, \ldots, F_F$ and outward normals~$n_k$
(not necessarily unit).
A \emph{transition} $F_i \to F_j$ is a point
$x \in F_i \cap F_j$ and $\varepsilon > 0$ such that
$x + t\, J_0 n_i \in K$ for $t \in [-\varepsilon, 0]$
and $x + t\, J_0 n_j \in K$ for $t \in [0, \varepsilon]$.
A transition $F_i \to F_j$ exists if and only if
\begin{enumerate}
\item $\omega_0(n_i, n_j) \ge 0$, and
\item $\exists\, x \in F_i \cap F_j$ with
  $\langle n_k, x \rangle < h_k$
  for all $k \notin \{i,j\}$ satisfying
  $\omega_0(n_i, n_k) < 0$ or
  $\omega_0(n_j, n_k) > 0$.
\end{enumerate}
\end{lemma}

\begin{proof}
Throughout, we use
$\langle n_k, J_0 n_\ell \rangle
= \omega_0(n_\ell, n_k)$.
Call $k \notin \{i,j\}$ \emph{blocking} if
$\omega_0(n_i, n_k) < 0$ or $\omega_0(n_j, n_k) > 0$.

\emph{Necessity.}
Let $x \in F_i \cap F_j$ realize the transition.
For the approach ($t < 0$) to satisfy
the~$F_j$ constraint:
$\langle n_j, x + t\, J_0 n_i \rangle
= h_j + t\,\omega_0(n_i, n_j) \leq h_j$,
which forces $\omega_0(n_i, n_j) \geq 0$
(since $t < 0$).
The departure gives the same by
antisymmetry of~$\omega_0$.
This proves~(1).

For~(2), suppose $k$ is blocking with
$\omega_0(n_i, n_k) < 0$.
Then
$\langle n_k, x + t\, J_0 n_i \rangle
= \langle n_k, x \rangle + t\,\omega_0(n_i, n_k)$.
Since both factors are negative, the product is
positive, so
$\langle n_k, x + t\, J_0 n_i \rangle
> \langle n_k, x \rangle$.
If $\langle n_k, x \rangle = h_k$, this
exceeds~$h_k$, violating $x + t\, J_0 n_i \in K$.
Hence $\langle n_k, x \rangle < h_k$.
The case $\omega_0(n_j, n_k) > 0$ is analogous
for the departure.

\emph{Sufficiency.}
Given $x \in F_i \cap F_j$ satisfying~(1) and~(2),
let $\delta$ be the minimum of
$h_k - \langle n_k, x \rangle > 0$
over all blocking~$k$.
For non-blocking $k \notin \{i,j\}$:
$\omega_0(n_i, n_k) \geq 0$, so
$\langle n_k, x + t\, J_0 n_i \rangle
\leq \langle n_k, x \rangle \leq h_k$
for $t \leq 0$;
similarly $\omega_0(n_j, n_k) \leq 0$, so
$\langle n_k, x + t\, J_0 n_j \rangle
\leq \langle n_k, x \rangle \leq h_k$
for $t \geq 0$.
For blocking~$k$, choosing
$\varepsilon < \delta / \max_k(
|\omega_0(n_i, n_k)| + |\omega_0(n_j, n_k)|)$
ensures all constraints hold.
\end{proof}

\begin{remark}[Numerical pruning of transitions]%
\label{rem:numerical-pruning}
The transition feasibility conditions in
Lemma~\ref{lem:transition-feasibility-appendix}
are evaluated as follows:
\begin{itemize}
\item \emph{Condition~(1): $\omega_0(n_i, n_j) \geq 0$.}
  This is decided exactly from the symplectic sign
  matrix~$D$ (Remark~\ref{rem:exact-quantities}).
  If $D(i, j) = {-}$, no transition
  $F_i \to F_j$ exists and the pair is pruned.
\item \emph{Condition~(2): LP feasibility.}
  This is a linear feasibility problem on
  the face $F_i \cap F_j$,
  solved numerically in \texttt{f64}.
  Following the three-valued framework of
  Section~\ref{app:general-approach},
  the numerical solver returns
  $\textsc{true}$ (feasible),
  $\textsc{unknown}$ (inconclusive), or
  $\textsc{false}$ (infeasible).
  Pairs with $\textsc{false}$ are pruned;
  pairs with $\textsc{true}$ or $\textsc{unknown}$
  are retained in the search space.
\end{itemize}
This combines exact and numerical pruning:
$D(i, j) = {-}$ is an exact, irrevocable prune;
the LP verdict is a sound but potentially
conservative three-valued prune.
\end{remark}

\begin{definition}[Simple polytope]%
\label{def:simple-polytope}
A polytope $K \subset \mathbb{R}^4$ is \emph{simple}
if every vertex lies on exactly $4$~facets,
i.e.\ every row of the incidence matrix~$E$
has exactly $4$ ones.
\end{definition}

\begin{remark}[Scope: simple polytopes]%
\label{rem:simple-scope}
We restrict to simple polytopes throughout.
This covers all polytopes arising in our experiments:
random polytopes are almost surely simple,
Lagrangian products are always simple,
and polytopes from future optimisation procedures
(e.g.\ gradient descent on the systolic ratio)
can be perturbed to simple ones if needed.
\end{remark}

\begin{definition}[Numerically robust polytope]%
\label{def:numerically-robust}
A simple polytope $K = K(n, h)$ with rational
coordinates is \emph{numerically robust}
if the following quantities, computed exactly
over~$\mathbb{Q}$, are bounded away from zero:
\begin{enumerate}
\item \emph{Gap margin:}
  $\gamma_{\min} \coloneqq
  \min_{j,\, i \notin S_j}
  (h_i - \langle n_i, v_j \rangle) > 0$,
  where $S_j$ is the set of facets through
  vertex~$v_j$.
\item \emph{Determinant margin:}
  $\delta_{\min} \coloneqq
  \min_j |\!\det(N_{S_j})| > 0$.
\item \emph{Symplectic margin} (when applicable):
  $\omega_{\min} \coloneqq
  \min_{\substack{i < k \\[1pt]
  D(i,k) \neq 0}}
  |\omega_0(n_i, n_k)| > 0$.
\end{enumerate}
These margins quantify how far each exact predicate
is from its decision boundary.
When all margins are large relative to
\texttt{f64} rounding error,
the rounded \texttt{f64} polytope preserves:
the vertex--facet incidence (via $\gamma_{\min}$
and $\delta_{\min}$),
and the nonzero entries of~$D$
(via $\omega_{\min}$).
% [GAP - AGENT CONFIDENCE 85%: The preservation claim is a standard
%  perturbation argument. No proof provided.
%  JÖRN: decide whether a proof or citation is needed.]
\end{definition}

\begin{remark}[Conversion to \texttt{f64}]%
\label{rem:f64-conversion}
A rational polytope $K = K(n, h)$ is converted to
\texttt{f64} by normalizing each normal to a unit vector:
\[
  \hat n_i = \frac{n_i}{\|n_i\|}, \qquad
  \hat h_i = \frac{h_i}{\|n_i\|}.
\]
The halfspace $\langle n_i, x \rangle \leq h_i$ is
equivalent to
$\langle \hat n_i, x \rangle \leq \hat h_i$.
The normalization and the rational-to-\texttt{f64}
conversion each introduce $O(\varepsilon_{\mathrm{mach}})$
rounding error per coordinate.
The capacity algorithm uses:
\begin{itemize}
\item the \emph{exact} rational data for discrete
  decisions (which facet subsets to enumerate,
  the symplectic sign matrix~$D$,
  and the exact part of transition pruning),
\item the \emph{rounded} \texttt{f64} coordinates
  for continuous computation (the SVD-based
  KKT solver, and the LP feasibility check
  in Remark~\ref{rem:numerical-pruning}).
\end{itemize}
\end{remark}
